\chapter{Linear Programming}\label{ch:lp}

\begin{center}\itshape
A linear function on a polyhedron takes its maximum at a corner;\\
walk from corner to better corner, and when no edge climbs, a certificate of optimality is waiting.\\
\upshape(the simplex method and duality, in our words)
\end{center}

In 1939 a graduate student at Berkeley arrived late to a lecture of the statistician Jerzy Neyman, copied the two problems written on the blackboard, and took them home as
homework. They seemed harder than usual. A few weeks later he handed in his solutions, apologising for the delay; the problems had been two famous open problems of statistics,
and the solutions became his thesis. The student was George Dantzig, and the story, which is so good that it is often retold as an urban legend, is true: it is told by
Dantzig himself and by his colleagues Richard Cottle and Margaret Wright.

During the war Dantzig worked for the US Air Force on what was then called \emph{programming}: planning schedules of training, supply and deployment. By mid-1947 he had
formulated such planning problems as the maximisation of a linear function subject to linear constraints, and by the end of that summer he had found a method to solve
them, the \emph{simplex method}. In October 1947 he explained his ideas to John von Neumann, who answered with an impromptu lecture on what we now call duality, which he
recognised from his own theory of games. (Dantzig told this meeting in his 2002 reminiscences; we paraphrase.) The simplex method went on to become one of the most used
algorithms in the world, and it is still taught exactly as Dantzig designed it, with one small rule added by Robert Bland in 1977 to make sure it always stops.

This chapter proves everything: that the simplex method stops (Bland's rule), that the optimum, when it exists, is reached at a corner, and the duality theorem, which
we read off the final step of the algorithm. Duality gives the Farkas lemma, the linear analogue of the theorem of the alternative of Tome~S1, Chapter~1. We then explain
why so many combinatorial linear programs have integral optimal solutions (total unimodularity), recover K\H onig's theorem and the integrality of flows, and finally keep
a promise of Tome~C1, Chapter~10, by proving Edmonds' matroid intersection theorem with linear programming.

\begin{think}
Maximise $5x_1+4x_2$ over the points of the plane with $x_1,x_2\ge0$, $x_1\le4$ and $x_2\le2$. Why must an optimal point be a corner of the rectangle? Can an optimum also be
attained at a point which is not a corner? Can there be a linear objective with no optimal corner at all, on a region bounded by finitely many lines?
\end{think}

\section{Linear programs and the bakery}\label{sec:lpintro}

A small bakery makes two kinds of trays: $x_1$ trays of croissants and $x_2$ trays of brioches each morning. A tray of croissants needs $6$\,kg of flour and $1$\,kg of butter, a tray of
brioches $4$\,kg of flour and $2$\,kg of butter; there are $24$\,kg of flour and $6$\,kg of butter. The baker never makes more than one tray of brioches beyond the croissants
($x_2\le x_1+1$), and never more than $2$ trays of brioches. A tray of croissants brings $5$ euros, a tray of brioches $4$. The best morning is the solution of
\begin{equation}\label{eq:bakery}
\begin{aligned}
\text{maximise}\quad & 5x_1+4x_2\\
\text{subject to}\quad & 6x_1+4x_2\le24,\quad x_1+2x_2\le6,\\
& -x_1+x_2\le1,\quad x_2\le2,\quad x_1,x_2\ge0.
\end{aligned}
\end{equation}
Figure~\ref{fig:bakery} shows the region of feasible mornings, a hexagon, and the lines on which the profit is constant, all perpendicular to $c=(5,4)$. The best morning is the
last point of the hexagon met by these lines as the profit grows: the corner $(3,\tfrac32)$, with profit $21$.

\begin{figure}[ht]
\centering
\fig{ch07-bakery}
\caption{The bakery's feasible mornings. The dotted line is the level set $5x_1+4x_2=21$; the simplex method walks along the red edges from $(0,0)$ to the optimal corner $(3,\tfrac32)$.}
\label{fig:bakery}
\end{figure}

In general, a linear program in \emph{standard form} is
\begin{equation}\label{eq:primal}
\text{maximise } c^{\mathsf T}x\quad\text{subject to}\quad Ax\le b,\ x\ge0,
\end{equation}
for $A\in\mathbb R^{m\times n}$, $b\in\mathbb R^m$ and $c\in\mathbb R^n$, the inequalities between vectors being taken coordinatewise. A point satisfying the constraints is
\emph{feasible}; the program is \emph{infeasible} if there is none, \emph{unbounded} if $c^{\mathsf T}x$ takes arbitrarily large values on feasible points, and a feasible $x$ is
\emph{optimal} if it maximises $c^{\mathsf T}x$. Every linear program can be put in standard form: minimising $c^{\mathsf T}x$ is maximising $-c^{\mathsf T}x$; a constraint $a^{\mathsf T}x\ge\beta$ is
$-a^{\mathsf T}x\le-\beta$; an equation is two inequalities; and a variable $x_j$ of any sign is written $x_j^+-x_j^-$ with $x_j^\pm\ge0$. So the three possible outcomes we
shall prove for~\eqref{eq:primal} hold for all linear programs.

How does the baker know that $21$ is the best? Add $\tfrac34$ times the flour constraint to $\tfrac12$ times the butter constraint: for every feasible morning,
\[
5x_1+4x_2\le\bigl(\tfrac34\cdot6+\tfrac12\cdot1\bigr)x_1+\bigl(\tfrac34\cdot4+\tfrac12\cdot2\bigr)x_2\le\tfrac34\cdot24+\tfrac12\cdot6=21.
\]
The multipliers form a certificate of optimality. Looking for the best certificate is itself a linear program, the \emph{dual} of~\eqref{eq:primal}:
\begin{equation}\label{eq:dual}
\text{minimise } b^{\mathsf T}y\quad\text{subject to}\quad A^{\mathsf T}y\ge c,\ y\ge0.
\end{equation}

\begin{proposition}[weak duality]\label{prop:weakduality}
If $x$ is feasible for~\eqref{eq:primal} and $y$ for~\eqref{eq:dual}, then $c^{\mathsf T}x\le b^{\mathsf T}y$. If equality holds, $x$ and $y$ are both optimal.
\end{proposition}
\begin{proof}
Since $x\ge0$ and $A^{\mathsf T}y\ge c$, and then $y\ge0$ and $Ax\le b$: $c^{\mathsf T}x\le(A^{\mathsf T}y)^{\mathsf T}x=y^{\mathsf T}(Ax)\le y^{\mathsf T}b$. If equality holds, any feasible $x'$ has $c^{\mathsf T}x'\le b^{\mathsf T}y=c^{\mathsf T}x$, and symmetrically for $y$.
\end{proof}

In particular, if the primal is unbounded the dual is infeasible, and conversely. The duality theorem will say that the best certificate always closes the gap.

\section{The simplex method}\label{sec:simplex}

We follow the presentation by \emph{dictionaries} of Va\v sek Chv\'atal's textbook. Assume first that $b\ge0$, so that $x=0$ is feasible. Introduce \emph{slack variables}
$x_{n+i}=b_i-\sum_{j=1}^na_{ij}x_j$ for $i=1,\dots,m$. The constraints become: all $n+m$ variables are nonnegative, and they satisfy the $m$ linear equations defining the slacks. Call
$\mathcal S\subseteq\mathbb R^{n+m}$ the set of solutions of these equations (signs ignored); it is an affine subspace of dimension $n$, on which $x_1,\dots,x_n$ are free and the slacks determined.
The objective is the function $z=\sum_jc_jx_j$ on $\mathcal S$.

\begin{definition}
A \emph{dictionary} is a partition of the indices $\{1,\dots,n+m\}$ into a set $B$ of $m$ \emph{basic} and a set $N$ of $n$ \emph{nonbasic} indices, together with equations
\[
x_i=\beta_i-\sum_{j\in N}\alpha_{ij}x_j\quad(i\in B),\qquad z=\zeta+\sum_{j\in N}\gamma_jx_j,
\]
which are satisfied by every point of $\mathcal S$ and define $\mathcal S$: a choice of the nonbasic variables, completed by the formulas, is a point of $\mathcal S$. Its \emph{basic solution} sets
$x_N=0$ and $x_B=\beta$; the dictionary is \emph{feasible} if $\beta\ge0$, and then its basic solution is a feasible point of value $\zeta$.
\end{definition}

The definitions of the slacks form the first dictionary, with $B=\{n+1,\dots,n+m\}$, $\beta=b$ and $\gamma=c$; it is feasible because $b\ge0$. For the bakery:
\[
\begin{aligned}
x_3&=24-6x_1-4x_2, & x_4&=6-x_1-2x_2,\\
x_5&=1+x_1-x_2, & x_6&=2-x_2,\\
z&=5x_1+4x_2.
\end{aligned}
\]

\begin{lemma}\label{lem:dictunique}
A dictionary is determined by its set of basic indices.
\end{lemma}
\begin{proof}
Take two dictionaries with the same $B$ and subtract their equations for $x_i$, $i\in B$: $\beta_i-\beta'_i=\sum_{j\in N}(\alpha_{ij}-\alpha'_{ij})x_j$ for every point of $\mathcal S$, that is, for every choice of $x_N$.
Taking $x_N=0$ gives $\beta_i=\beta'_i$, and taking for $x_N$ the unit vectors gives $\alpha_{ij}=\alpha'_{ij}$. The same argument applies to the row of $z$, which is a fixed function on~$\mathcal S$.
\end{proof}

A \emph{pivot} exchanges one basic and one nonbasic variable. Given a feasible dictionary, choose an \emph{entering} index $e\in N$ with $\gamma_e>0$: increasing $x_e$ from $0$, the other nonbasic
variables staying at $0$, increases $z$. The basic variables change as $x_i=\beta_i-\alpha_{ie}x_e$; those with $\alpha_{ie}>0$ decrease, and the first to reach $0$ stops the increase. So choose a
\emph{leaving} index $\ell\in B$ with $\alpha_{\ell e}>0$ minimising the ratio $\beta_\ell/\alpha_{\ell e}$ (the \emph{ratio test}), solve the equation of $x_\ell$ for $x_e$, and substitute it in all the other equations.

\begin{proposition}\label{prop:pivot}
Let a feasible dictionary be given.
\begin{enumerate}
\item If $\gamma_j\le0$ for all $j\in N$, its basic solution is optimal.
\item If $\gamma_e>0$ and $\alpha_{ie}\le0$ for all $i\in B$, the program is unbounded.
\item Otherwise a pivot with the ratio test gives a feasible dictionary, with basis $B-\ell+e$, and $\zeta$ increases by $\gamma_e\beta_\ell/\alpha_{\ell e}\ge0$. The basic solution does not change if $\beta_\ell=0$.
\end{enumerate}
\end{proposition}
\begin{proof}
(1) Every feasible point lies in $\mathcal S$ and has $x_N\ge0$, so $z=\zeta+\sum\gamma_jx_j\le\zeta$. (2) The points $x_e=t$, $x_{N-e}=0$, $x_i=\beta_i-\alpha_{ie}t$ are feasible for all $t\ge0$, and
$z=\zeta+\gamma_et\to\infty$. (3) Write $t=\beta_\ell/\alpha_{\ell e}\ge0$. The new equations hold on $\mathcal S$ because they are obtained from the old ones by solving and substituting, and they define $\mathcal S$ because the
substitution can be undone; so they form a dictionary. Its basic solution is the point $x_e=t$, $x_{N-e}=0$, $x_i=\beta_i-\alpha_{ie}t$ (with $x_\ell=0$): indeed in the new dictionary $x_\ell$ is nonbasic, and this point has all new
nonbasic variables equal to $0$. It is nonnegative: $\beta_i-\alpha_{ie}t\ge\beta_i\ge0$ if $\alpha_{ie}\le0$, and $\beta_i-\alpha_{ie}t\ge0$ by the choice of $\ell$ otherwise. Its value is $\zeta+\gamma_et$. If $t=0$ it is the old basic solution.
\end{proof}

For the bakery, $x_1$ enters; the ratios are $24/6=4$ for $x_3$ and $6/1$ for $x_4$ ($x_5$ and $x_6$ do not decrease), so $x_3$ leaves, and $x_1=4-\tfrac23x_2-\tfrac16x_3$. Substituting:
\[
\begin{aligned}
x_1&=4-\tfrac23x_2-\tfrac16x_3, & x_4&=2-\tfrac43x_2+\tfrac16x_3,\\
x_5&=5-\tfrac53x_2-\tfrac16x_3, & x_6&=2-x_2,\\
z&=20+\tfrac23x_2-\tfrac56x_3.
\end{aligned}
\]
Now $x_2$ enters; the ratios are $6$, $\tfrac32$, $3$, $2$, so $x_4$ leaves, and after substitution
\[
\begin{aligned}
x_1&=3+\tfrac12x_4-\tfrac14x_3, & x_2&=\tfrac32-\tfrac34x_4+\tfrac18x_3,\\
x_5&=\tfrac52+\tfrac54x_4-\tfrac38x_3, & x_6&=\tfrac12+\tfrac34x_4-\tfrac18x_3,\\
z&=21-\tfrac34x_3-\tfrac12x_4.
\end{aligned}
\]
All coefficients of $z$ are negative: the morning $(3,\tfrac32)$ is optimal, and the dictionary itself proves it. This is the red path of Figure~\ref{fig:bakery}.

When $\beta_\ell=0$ the pivot is \emph{degenerate}: the dictionary changes but not the point. A sequence of degenerate pivots may come back to a dictionary already seen, and then a deterministic method
loops forever: this is \emph{cycling}. It cannot happen without degeneracy (Exercise~\ref{ex:degenerate}), but with it, it does: Alan Hoffman in 1953 and Martin Beale in 1955 built examples on which natural rules cycle, and the notebook runs one. Bland's rule removes it.

\begin{theorem}[Bland, 1977]\label{thm:bland}
Choose as entering variable the one of smallest index among those with $\gamma_j>0$, and as leaving variable the one of smallest index among those that tie in the ratio test. Then the simplex method never
cycles; it stops after at most $\binom{n+m}{m}$ pivots, with an optimal dictionary or a proof of unboundedness.
\end{theorem}
\begin{proof}
Suppose that the dictionaries $D_0,D_1,\dots,D_k=D_0$ form a cycle. The value $\zeta$ never decreases (Proposition~\ref{prop:pivot}) and comes back to its starting value, so every pivot of the cycle is degenerate, and
all these dictionaries have the same basic solution $\bar x$. Call an index \emph{fickle} if it is basic in some $D_i$ and nonbasic in another. A fickle variable is $0$ in the basic solution of a dictionary where it is nonbasic,
so $\bar x_j=0$ for every fickle $j$.

Let $t$ be the largest fickle index. Let $D$ be a dictionary of the cycle in which $x_t$ leaves, say with $x_s$ entering, and let $D^*$ be one in which $x_t$ enters; $s$ is fickle, so $s<t$. Write $D$ as above with
coefficients $\beta,\alpha,\gamma$ and basis $B$, and write the objective row of $D^*$ as $z=\zeta+\sum_{j=1}^{n+m}\gamma^*_jx_j$, with $\gamma^*_j=0$ for $j$ basic in $D^*$ (the constant is the same $\zeta$).

For every real $u$, the point of $\mathcal S$ with $x_s=u$, the other variables of $N$ at $0$, and $x_i=\beta_i-\alpha_{is}u$ for $i\in B$, has $z=\zeta+\gamma_su$ by $D$. The equation of $D^*$ also holds there:
\[
\zeta+\gamma_su=\zeta+\gamma^*_su+\sum_{i\in B}\gamma^*_i(\beta_i-\alpha_{is}u).
\]
Both sides are affine in $u$; comparing the coefficients of $u$,
\begin{equation}\label{eq:blandkey}
\gamma_s-\gamma^*_s=-\sum_{i\in B}\gamma^*_i\alpha_{is}.
\end{equation}
Now $\gamma_s>0$, since $x_s$ enters in $D$. And $\gamma^*_s\le0$: otherwise, as $s<t$, Bland's rule would have chosen $x_s$ rather than $x_t$ to enter in $D^*$. So the left side of~\eqref{eq:blandkey} is positive, and
some $r\in B$ has $\gamma^*_r\alpha_{rs}<0$. Then $\gamma^*_r\ne0$, so $r$ is nonbasic in $D^*$ while basic in $D$: it is fickle, and $r\le t$.

If $r=t$: $\alpha_{ts}>0$, since $x_t$ leaves in $D$ when $x_s$ enters, and $\gamma^*_t>0$, since $x_t$ enters in $D^*$; so $\gamma^*_t\alpha_{ts}>0$, which is false. Hence $r<t$. Then $\gamma^*_r\le0$, because $r<t$ and $x_t$ was chosen
in $D^*$; as $\gamma^*_r\ne0$, $\gamma^*_r<0$, and therefore $\alpha_{rs}>0$. But $r$ is fickle, so $\beta_r=\bar x_r=0$, and $x_r$ was a candidate in the ratio test of $D$ with ratio $0$, the smallest possible, tying with $x_t$
(whose ratio is $\beta_t/\alpha_{ts}=0$ as well). Bland's rule would have made $x_r$ leave, not $x_t$, since $r<t$. This contradiction shows that there is no cycle.

So no dictionary is visited twice. A dictionary is determined by its basis (Lemma~\ref{lem:dictunique}), there are at most $\binom{n+m}{m}$ bases, and the method stops, in case~(1) or~(2) of Proposition~\ref{prop:pivot}.
\end{proof}

It remains to find a first feasible dictionary when $b\not\ge0$. Dantzig's answer is to solve another linear program first.

\begin{theorem}[the fundamental theorem of linear programming]\label{thm:fundamental}
A linear program is infeasible, or unbounded, or has an optimal solution. For~\eqref{eq:primal}, the two-phase simplex method with Bland's rule decides which, and in the third case ends with an optimal dictionary.
\end{theorem}
\begin{proof}
By the reduction of Section~\ref{sec:lpintro}, it is enough to treat~\eqref{eq:primal}, and if $b\ge0$ Theorem~\ref{thm:bland} does it. Otherwise, introduce a new variable $x_0$ and the \emph{auxiliary} program
\[
\text{maximise } -x_0\quad\text{subject to}\quad \sum_ja_{ij}x_j-x_0\le b_i\ (1\le i\le m),\quad x_0,x_1,\dots,x_n\ge0.
\]
It is feasible ($x=0$ and $x_0$ large), and~\eqref{eq:primal} is feasible exactly when the auxiliary program has a feasible point with $x_0=0$, that is, when its optimum is $0$. Its first dictionary
$x_{n+i}=b_i-\sum_ja_{ij}x_j+x_0$ is not feasible, but one pivot repairs it: let $x_0$ enter and let leave the slack $x_{n+k}$ with $b_k$ most negative. Then $x_0=-b_k+\sum_ja_{kj}x_j+x_{n+k}$, and every
$x_{n+i}=b_i-b_k+\dots$ has a nonnegative constant. The objective $-x_0$ is bounded above by $0$, so the simplex method with Bland's rule (phase one) ends with an optimal dictionary of the auxiliary program.

If its value is negative, \eqref{eq:primal} is infeasible. If it is $0$, the basic solution has $x_0=0$. If $x_0$ is basic, its row reads $x_0=0-\sum_{j\in N}\alpha_{0j}x_j$, and some $\alpha_{0j}\ne0$: otherwise $x_0$ would be $0$ at every
solution of the auxiliary equations, whereas $x_0$ can be chosen freely there. A pivot exchanging $x_0$ and $x_j$ has $\beta_0=0$, so it keeps the basic solution and the dictionary feasible, whatever the sign of $\alpha_{0j}$. Now $x_0$ is nonbasic. Delete it from all
equations: setting $x_0=0$ in a dictionary of the auxiliary program gives a dictionary for the slack equations of~\eqref{eq:primal}, which is feasible. Finally replace the row of $z$ by $c^{\mathsf T}x$ written in terms of the nonbasic
variables, by substituting the rows of the basic $x_j$, and run Theorem~\ref{thm:bland} (phase two).
\end{proof}

\begin{think}
In phase one, why may $x_0$ end up basic at value $0$, rather than leaving the basis? Find a small example where the auxiliary optimum is $0$ but phase one, pivoting carelessly, keeps $x_0$ in the basis.
\end{think}

\section{Duality}\label{sec:duality}

The final dictionary of the bakery says $z=21-\tfrac34x_3-\tfrac12x_4$, and the certificate found by hand in Section~\ref{sec:lpintro} was $y=(\tfrac34,\tfrac12,0,0)$: the dual solution is written in the objective row, under the slack
variables. This is no accident.

\begin{theorem}[strong duality; von Neumann, 1947; Gale, Kuhn and Tucker, 1951]\label{thm:strongduality}
If the primal~\eqref{eq:primal} has an optimal solution, so has the dual~\eqref{eq:dual}, and the optimal values are equal. More precisely, if $z=\zeta+\sum_{j\in N}\gamma_jx_j$ is the objective row of an optimal dictionary,
then $y_i=-\gamma_{n+i}$ (with $\gamma_k=0$ for basic $k$) is an optimal dual solution.
\end{theorem}
\begin{proof}
By Theorem~\ref{thm:fundamental} there is an optimal dictionary, with $\gamma\le0$ and $\zeta$ the optimal value. Its objective row holds at every point of $\mathcal S$: for every $x\in\mathbb R^n$, with the slacks
$x_{n+i}=b_i-\sum_ja_{ij}x_j$,
\[
\sum_{j=1}^nc_jx_j=\zeta+\sum_{j=1}^n\gamma_jx_j-\sum_{i=1}^my_i\Bigl(b_i-\sum_{j=1}^na_{ij}x_j\Bigr).
\]
Both sides are affine functions of $x\in\mathbb R^n$; comparing constants gives $\zeta=\sum_iy_ib_i=b^{\mathsf T}y$, and comparing the coefficients of $x_j$ gives $\sum_ia_{ij}y_i=c_j-\gamma_j\ge c_j$. So $y=-(\gamma_{n+1},\dots,\gamma_{n+m})\ge0$ is dual
feasible with $b^{\mathsf T}y=\zeta$, and both are optimal by Proposition~\ref{prop:weakduality}.
\end{proof}

The dual of~\eqref{eq:dual}, put in standard form, is~\eqref{eq:primal} again (Exercise~\ref{ex:dualrules}), so the theorem also holds with the roles exchanged. Together with Proposition~\ref{prop:weakduality} and
Theorem~\ref{thm:fundamental}, it leaves four possibilities: both programs optimal with equal values; primal unbounded and dual infeasible; dual unbounded and primal infeasible; or both infeasible.
The economists read the dual variables as \emph{prices}: $y_1=\tfrac34$ is what one more kilogram of flour is worth to the bakery (Exercise~\ref{ex:bakerydual}).

\begin{corollary}[complementary slackness]\label{cor:compslack}
Feasible $x$ and $y$ are both optimal if and only if $y_i=0$ whenever the $i$th constraint is not tight at $x$, and $x_j=0$ whenever the $j$th dual constraint is not tight at $y$.
\end{corollary}
\begin{proof}
In the proof of Proposition~\ref{prop:weakduality}, $b^{\mathsf T}y-c^{\mathsf T}x=y^{\mathsf T}(b-Ax)+(A^{\mathsf T}y-c)^{\mathsf T}x$ is a sum of products of nonnegative numbers. It vanishes exactly under the stated conditions,
and by Theorem~\ref{thm:strongduality} it vanishes exactly when $x$ and $y$ are optimal.
\end{proof}

Tome~S1, Corollary~1.14, showed that a system of linear equations with no solution has a short proof of that fact: a combination of the equations reading $0=1$. Inequalities have their own certificate, found by the
Hungarian physicist Gyula Farkas in 1902.

\begin{corollary}[Farkas' lemma]\label{cor:farkas}
The system $Ax\le b$, $x\ge0$ has a solution if and only if $y^{\mathsf T}b\ge0$ for every $y\ge0$ with $A^{\mathsf T}y\ge0$. Equivalently, exactly one of the following holds: the system has a solution, or there is
$y\ge0$ with $A^{\mathsf T}y\ge0$ and $y^{\mathsf T}b<0$.
\end{corollary}
\begin{proof}
If $x$ is a solution and $y\ge0$, $A^{\mathsf T}y\ge0$, then $y^{\mathsf T}b\ge y^{\mathsf T}Ax=(A^{\mathsf T}y)^{\mathsf T}x\ge0$. Conversely, suppose there is no solution, and consider the program
$\text{maximise }-b^{\mathsf T}y$ subject to $-A^{\mathsf T}y\le0$, $y\ge0$. It is feasible ($y=0$). If it had an optimal solution, Theorem~\ref{thm:strongduality} would give an optimal solution of its dual, which is:
minimise $0$ subject to $-Ax\ge-b$, $x\ge0$, a solution of our system. So it has no optimal solution, and by Theorem~\ref{thm:fundamental} it is unbounded: some feasible $y$ has $-b^{\mathsf T}y>0$.
\end{proof}

Geometrically, if $b$ cannot be reached, a hyperplane separates it from what can: the certificate $y$ is its normal vector. Farkas' lemma, the strong duality theorem and von Neumann's minimax theorem for two-person
zero-sum games (Exercise~\ref{ex:gamevalue}) are three faces of one fact, and each can be deduced from the others.

\section{Vertices}\label{sec:vertices}

The simplex method walks on corners. We now describe the corners of any polyhedron directly; this is what makes integrality questions tractable.

\begin{definition}
Let $P=\{x\in\mathbb R^n:Ax\le b\}$, with rows $a_1^{\mathsf T},\dots,a_m^{\mathsf T}$ of $A$. A constraint $i$ is \emph{tight} at $x\in P$ if $a_i^{\mathsf T}x=b_i$. A point $x\in P$ is a \emph{vertex} of $P$ if the rows $a_i$
of the constraints tight at $x$ span $\mathbb R^n$.
\end{definition}

A vertex is the unique solution of the equations of its tight constraints, and there are finitely many sets of constraints, so a polyhedron has finitely many vertices. Exercise~\ref{ex:extreme} shows that vertices are the
\emph{extreme points}, those that are not the midpoint of two other points of $P$, and Exercise~\ref{ex:bfsvertex} that, for the standard form, they are the basic solutions of the feasible dictionaries.

\begin{theorem}\label{thm:vertex}
Let $P=\{x:Ax\le b\}$ be nonempty, with $A$ of rank $n$. Then $P$ has a vertex, and if $c^{\mathsf T}x$ is bounded above on $P$, its maximum on $P$ exists and is attained at a vertex.
\end{theorem}
\begin{proof}
We show that from any $x\in P$ we can reach a vertex $v$ with $c^{\mathsf T}v\ge c^{\mathsf T}x$; this proves both claims (the first with $c=0$), since the maximum of $c$ over the finitely many vertices then bounds $c$ on all of $P$.

Let $T$ be the set of constraints tight at $x$, and suppose their rows span a subspace of dimension $k<n$. Pick $d\ne0$ orthogonal to all $a_i$, $i\in T$. Since $A$ has rank $n$, $Ad\ne0$. If $c^{\mathsf T}d\ne0$, choose the
sign of $d$ so that $c^{\mathsf T}d>0$; then some $a_i^{\mathsf T}d>0$, for otherwise $x+\lambda d\in P$ for all $\lambda\ge0$ and $c$ would be unbounded on $P$. If $c^{\mathsf T}d=0$, choose the sign of $d$ so that some $a_i^{\mathsf T}d>0$, which is
possible since $Ad\ne0$. In both cases $c^{\mathsf T}d\ge0$, and moving from $x$ to $x+\lambda d$ keeps the tight constraints tight; let $\lambda^*$ be the largest $\lambda$ with $x+\lambda d\in P$, namely the least of
$(b_i-a_i^{\mathsf T}x)/a_i^{\mathsf T}d$ over the $i$ with $a_i^{\mathsf T}d>0$ (none of which is in $T$). At $x'=x+\lambda^*d$ a new constraint $j$ is tight, and $a_j^{\mathsf T}d>0$ while all $a_i$, $i\in T$, are orthogonal to $d$, so $a_j$ is not in their span.
Hence the tight rows at $x'$ span a space of dimension at least $k+1$, and $c^{\mathsf T}x'\ge c^{\mathsf T}x$. After at most $n$ steps we reach a vertex.
\end{proof}

Polyhedra containing a line, such as a half-plane, have no vertex; requiring $x\ge0$, as in the standard form, puts the rows of the identity among the constraints and excludes them.

\section{Total unimodularity}\label{sec:tu}

Many combinatorial problems are integer programs: maximise $c^{\mathsf T}x$ over the \emph{integer} points of a polyhedron. In general this is much harder (Tome~C3 will show that it is NP-hard). But sometimes the
polyhedron has only integral vertices, and then the linear program, which we can solve, has an integral optimal solution, which is the answer.

\begin{definition}
An integer matrix is \emph{totally unimodular} (TU) if every square submatrix has determinant $0$, $1$ or $-1$.
\end{definition}

In particular the entries of a TU matrix are in $\{0,\pm1\}$. Transposing a TU matrix, deleting rows or columns, multiplying a row by $-1$, repeating a row, or adding a row or column with a single nonzero entry $\pm1$
keeps it TU: a square submatrix of the new matrix is, up to sign, a square submatrix of the old one, or has two equal or opposite rows (determinant $0$), or can be expanded along the added line.

\begin{theorem}[Hoffman and Kruskal, 1956]\label{thm:hoffmankruskal}
If $A$ is TU and $b$ is integral, every vertex of $P=\{x:Ax\le b\}$ is integral. Consequently, if moreover $P$ is nonempty, $A$ has rank $n$ and $c^{\mathsf T}x$ is bounded above on $P$, its maximum on $P$ is attained at
an integral point.
\end{theorem}
\begin{proof}
At a vertex $v$, choose $n$ tight constraints with linearly independent rows: $A_0v=b_0$ with $A_0$ an invertible square submatrix of $A$, so $\det A_0=\pm1$. By Cramer's rule each coordinate of $v=A_0^{-1}b_0$ is a determinant of an integer
matrix divided by $\pm1$, an integer. The second sentence follows by Theorem~\ref{thm:vertex}.
\end{proof}

Hoffman and Kruskal also proved the converse: if $\{x\ge0:Ax\le b\}$ has integral vertices for every integral $b$, then $A$ is TU. Here is a criterion which covers the matrices of graphs.

\begin{theorem}[Heller and Tompkins, 1956]\label{thm:tucriterion}
Let $A$ have entries in $\{0,\pm1\}$, at most two nonzero entries in each column, and suppose its rows can be split into two groups $R_1$ and $R_2$ such that, in each column with two nonzero entries, these entries are in
different groups if they have the same sign, and in the same group if they have opposite signs. Then $A$ is TU.
\end{theorem}
\begin{proof}
A square submatrix $Q$ satisfies the same hypothesis, with the groups restricted to its rows. We show $\det Q\in\{0,\pm1\}$ by induction on its size; size $1$ is clear. If $Q$ has a zero column, $\det Q=0$. If a column has exactly one
nonzero entry, expanding along it gives $\pm$ the determinant of a smaller submatrix, which is $0$ or $\pm1$ by induction. Otherwise every column has exactly two nonzero entries; by hypothesis they cancel in
$\sum_{i\in R_1}(\text{row }i)-\sum_{i\in R_2}(\text{row }i)$, which is therefore $0$. This is a nontrivial vanishing combination of the rows, so $\det Q=0$.
\end{proof}

\begin{corollary}\label{cor:tugraphs}
The vertex--edge incidence matrix of a bipartite graph is TU. The vertex--arc incidence matrix of a directed graph, with $+1$ at the head and $-1$ at the tail of each arc, is TU.
\end{corollary}
\begin{proof}
For a bipartite graph take $R_1$ and $R_2$ to be the two sides: each column has two entries $1$, one on each side. For a digraph take all the rows in $R_1$: each column has a $+1$ and a $-1$.
\end{proof}

Here is what this gives. For a bipartite graph $G$ with incidence matrix $M$ (rows the vertices, columns the edges), the matching polytope $\{x\ge0:Mx\le\mathbf 1\}$ has integral vertices, which are $0/1$ vectors with at most
one edge at each vertex: matchings. So the linear program ``maximise $\mathbf 1^{\mathsf T}x$'' over it has value $\nu(G)$, the size of a largest matching. Its dual is ``minimise $\mathbf 1^{\mathsf T}y$ subject to $M^{\mathsf T}y\ge\mathbf 1$, $y\ge0$'': the matrix
of this system, $M^{\mathsf T}$ with sign changes and the rows of the identity, is TU, so the dual optimum is attained at an integral $y$. Lowering every entry larger than $1$ to $1$ keeps $y$ feasible and does not increase the objective, so there is an optimal
$0/1$ vector $y$, which is a set of vertices touching every edge: a vertex cover. Strong duality gives:

\begin{corollary}[K\H onig's theorem, by linear programming]\label{cor:koniglp}
In a bipartite graph the largest matching and the smallest vertex cover have the same size. In a network with integer capacities there is an integral maximum flow.
\end{corollary}
\begin{proof}
The first sentence was just proved. For flows, the conservation equations at the vertices other than $s$ and $t$ use rows of the incidence matrix of the digraph, written as two inequalities each, and the capacity constraints
$0\le f\le c$ add rows $\pm1$ of the identity: the system is TU with integral right side, and the value of the flow is a linear function, bounded by the total capacity. Theorem~\ref{thm:hoffmankruskal} applies.
\end{proof}

These were Tome~C1, Theorems~6.4 and~6.6; the point is that one theorem about determinants explains both, and that the proof of K\H onig's theorem came out of the duality of the linear programs. Exercise~\ref{ex:oddcycle} shows what goes wrong
for a triangle, and why Edmonds' blossoms of Chapter~\ref{ch:flows} were needed.

\begin{think}
The matrix of the \emph{consecutive ones} property, where each row is a block of consecutive $1$s, does not satisfy the hypothesis of Theorem~\ref{thm:tucriterion}. Is it TU? Try a few $3\times3$ examples before reading
Exercise~\ref{ex:tuinterval}.
\end{think}

\section{Matroid intersection by linear programming}\label{sec:matroidlp}

Let $M_1$ and $M_2$ be matroids on the same finite ground set $E$, with rank functions $r_1$ and $r_2$ (Tome~C1, Chapter~10). Recall that a rank function satisfies $0\le r(A)\le|A|$, is monotone, and is submodular:
$r(A\cup B)+r(A\cap B)\le r(A)+r(B)$ (Tome~C1, Theorem~10.6). Tome~C1 stated Edmonds' theorem
\[
\max\{|I|:I\text{ independent in }M_1\text{ and }M_2\}=\min_{A\subseteq E}\bigl(r_1(A)+r_2(E\setminus A)\bigr),
\]
and we now prove it. For $x\in\mathbb R^E$ and $A\subseteq E$ write $x(A)=\sum_{e\in A}x_e$, and $\chi^A$ for the incidence vector of $A$. The \emph{matroid polytope} of $M$ is
\[
P(M)=\{x\in\mathbb R^E:\ x\ge0,\ x(A)\le r(A)\text{ for all }A\subseteq E\},
\]
with one constraint for each of the $2^{|E|}$ subsets. Its integral points are exactly the vectors $\chi^I$, $I$ independent: such a point has $x_e\le r(\{e\})\le1$, so it is some $\chi^I$, and $|I|=x(I)\le r(I)$ makes $I$
independent; conversely $\chi^I(A)=|I\cap A|=r(I\cap A)\le r(A)$.

Fix $x\in P(M)$. Call $A$ \emph{tight} if $x(A)=r(A)$; the empty set is tight.

\begin{lemma}[uncrossing]\label{lem:uncross}
If $A$ and $B$ are tight, so are $A\cap B$ and $A\cup B$.
\end{lemma}
\begin{proof}
Using $x(A)+x(B)=x(A\cap B)+x(A\cup B)$, the constraints of $P(M)$ and submodularity,
\[
r(A)+r(B)=x(A\cap B)+x(A\cup B)\le r(A\cap B)+r(A\cup B)\le r(A)+r(B).
\]
So equality holds throughout, and since $x(A\cap B)\le r(A\cap B)$ and $x(A\cup B)\le r(A\cup B)$ separately, both are equalities.
\end{proof}

\begin{lemma}[chains span]\label{lem:chainspan}
Let $\mathcal F$ be a nonempty family of subsets of $E$ closed under union and intersection, and $\mathcal C$ a maximal chain in $\mathcal F$ (a subfamily totally ordered by inclusion, not contained in a larger one). Then
every $\chi^A$, $A\in\mathcal F$, is a linear combination of the $\chi^C$, $C\in\mathcal C$.
\end{lemma}
\begin{proof}
Write $\mathcal C=\{C_1\subsetneq C_2\subsetneq\dots\subsetneq C_k\}$. The intersection $L$ of all members of $\mathcal F$ is in $\mathcal F$ and contained in every set of the chain, so by maximality $L\in\mathcal C$, that is $L=C_1$; likewise the
union of all members is $C_k$. Let $A\in\mathcal F$ and $2\le i\le k$. The set $D=C_{i-1}\cup(A\cap C_i)$ is in $\mathcal F$ and satisfies $C_{i-1}\subseteq D\subseteq C_i$, so it is comparable with every member of the chain; by maximality
$D=C_{i-1}$ or $D=C_i$. In words: $A$ contains either none or all of $C_i\setminus C_{i-1}$. Since $C_1\subseteq A\subseteq C_k$, $A$ is $C_1$ together with some of the differences $C_i\setminus C_{i-1}$, and
$\chi^A=\chi^{C_1}+\sum_{i\in J}(\chi^{C_i}-\chi^{C_{i-1}})$ for some set $J$ of indices.
\end{proof}

\begin{lemma}[two chains]\label{lem:twochains}
Let $\mathcal C_1$ and $\mathcal C_2$ be chains of subsets of $E$. The matrix whose rows are the $\chi^C$, $C\in\mathcal C_1$, and the $\chi^C$, $C\in\mathcal C_2$, is TU.
\end{lemma}
\begin{proof}
A square submatrix $Q$ picks some sets of each chain and restricts them to a set $S$ of columns; the restrictions $C\cap S$ still form two chains. Order the rows coming from $\mathcal C_1$ as $D_1\subseteq D_2\subseteq\dots\subseteq D_p$ and, from the
last to the second, replace the row of $D_i$ by the row of $D_i$ minus the row of $D_{i-1}$; do the same with the rows from $\mathcal C_2$. Subtracting a row from another does not change the determinant. The new rows of each group are
incidence vectors of pairwise disjoint sets, so the new matrix has entries in $\{0,1\}$ and each column has at most one $1$ in each group. It satisfies Theorem~\ref{thm:tucriterion}, with the two groups as $R_1$ and $R_2$,
so its determinant, which is $\det Q$, is $0$ or $\pm1$.
\end{proof}

\begin{theorem}[Edmonds, 1970]\label{thm:matroidpolytope}
Every vertex of $P(M_1)\cap P(M_2)$ is integral, hence the incidence vector of a set independent in both matroids. In particular (taking $M_1=M_2$) the vertices of $P(M)$ are the incidence vectors of the independent sets.
\end{theorem}
\begin{proof}
Let $x$ be a vertex, $Z=\{e:x_e=0\}$, and $\mathcal F_k$ the family of sets tight for $M_k$ at $x$ ($k=1,2$). By Lemma~\ref{lem:uncross}, $\mathcal F_k$ is closed under union and intersection; it contains $\varnothing$. Let $\mathcal C_k$ be a maximal chain in
$\mathcal F_k$. The tight constraints at $x$ have rows $-\chi^{\{e\}}$ ($e\in Z$) and $\chi^A$ ($A\in\mathcal F_1\cup\mathcal F_2$), which span $\mathbb R^E$; by Lemma~\ref{lem:chainspan}, so do the rows $\chi^{\{e\}}$ ($e\in Z$) and $\chi^C$
($C\in\mathcal C_1\cup\mathcal C_2$). Hence $x$ is the unique solution of
\[
x_e=0\ (e\in Z),\qquad x(C)=r_1(C)\ (C\in\mathcal C_1),\qquad x(C)=r_2(C)\ (C\in\mathcal C_2).
\]
Let $S=E\setminus Z$. The vector $x_S$ is then the unique solution of the system $Nx_S=\rho$, where $N$ has rows $\chi^{C\cap S}$ for $C$ in the two chains and $\rho$ is the corresponding integral vector of ranks. So $N$ has rank $|S|$, and we can pick
$|S|$ linearly independent rows: an invertible square submatrix $N_0$ of $N$, TU by Lemma~\ref{lem:twochains}, with $\det N_0=\pm1$ and $N_0x_S=\rho_0$. By Cramer's rule $x_S$ is integral, and so is $x$. The last claim follows from the description of the
integral points of $P(M_k)$ above.
\end{proof}

\begin{theorem}[matroid intersection; Edmonds, 1970]\label{thm:matroidintersection}
For matroids $M_1$, $M_2$ on $E$,
\[
\max\{|I|:I\in\mathcal I_1\cap\mathcal I_2\}=\min_{A\subseteq E}\bigl(r_1(A)+r_2(E\setminus A)\bigr).
\]
\end{theorem}
\begin{proof}
The inequality $\le$ is easy: $|I|=|I\cap A|+|I\setminus A|\le r_1(A)+r_2(E\setminus A)$, as $I\cap A$ is independent in $M_1$ and $I\setminus A$ in $M_2$.

Let $\kappa$ be the left side. Consider the program: maximise $\mathbf 1^{\mathsf T}x$ over $P(M_1)\cap P(M_2)$. Its feasible region contains $0$, lies in $[0,1]^E$, and has the rows of the identity among its constraints, so by
Theorem~\ref{thm:vertex} its maximum is attained at a vertex, which is some $\chi^I$ with $I$ common independent by Theorem~\ref{thm:matroidpolytope}; since every such $\chi^I$ is feasible, the maximum is $\kappa$. The program is in
standard form, with a constraint $x(A)\le r_k(A)$ for each $k$ and each $A$. Its dual has variables $y_{k,A}\ge0$:
\begin{equation}\label{eq:mdual}
\begin{aligned}
&\text{minimise }\sum_{k,A}r_k(A)\,y_{k,A}\\
&\text{subject to }\sum_{k}\ \sum_{A\ni e}y_{k,A}\ge1\ \text{ for all }e\in E,\quad y\ge0,
\end{aligned}
\end{equation}
and by Theorem~\ref{thm:strongduality} its optimal value is $\kappa$.

\emph{Uncrossing the dual.} Let $\phi(A)=|A|\,(|E|-|A|)$. The program ``minimise $\sum_{k,A}\phi(A)\,y_{k,A}$ over the optimal solutions of~\eqref{eq:mdual}'' is a linear program (add the constraint that the objective of~\eqref{eq:mdual} is $\le\kappa$), feasible and
bounded below by $0$; by Theorem~\ref{thm:fundamental} it has an optimal solution $y$. Suppose that $y_{k,A}>0$ and $y_{k,B}>0$ for two \emph{crossing} sets, neither containing the other, and let $\varepsilon=\min(y_{k,A},y_{k,B})$. Decrease
$y_{k,A}$ and $y_{k,B}$ by $\varepsilon$ and increase $y_{k,A\cap B}$ and $y_{k,A\cup B}$ by $\varepsilon$. Since $\chi^A+\chi^B=\chi^{A\cap B}+\chi^{A\cup B}$, every element is covered as before, and by submodularity the objective of~\eqref{eq:mdual} does
not increase: the new $y$ is still optimal. But $a=|A\cap B|$ and $a'=|A\cup B|$ satisfy $a+a'=|A|+|B|$ with $a<\min(|A|,|B|)$ and $a'>\max(|A|,|B|)$, and $\phi(A)=t(|E|-t)$ with $t=|A|$ is a strictly concave function of $t$, so
$\phi(A\cap B)+\phi(A\cup B)<\phi(A)+\phi(B)$, contradicting the choice of $y$. Hence, for each $k$, the sets $A$ with $y_{k,A}>0$ form a chain $\mathcal C_k$.

\emph{Integrality.} Restrict~\eqref{eq:mdual} to the variables $y_{k,C}$, $C\in\mathcal C_k$. Its feasible points extend by zeros to feasible points of~\eqref{eq:mdual}, so its value is at least $\kappa$, and $y$ shows that it is $\kappa$. Its constraint
matrix is the transpose of the two-chain matrix of Lemma~\ref{lem:twochains}, with sign changes and the rows $-\chi^{\{(k,C)\}}$ of $y\ge0$ added: it is TU, and Theorem~\ref{thm:hoffmankruskal} (for $\max$ of the opposite objective, which is bounded above by $0$
since ranks are nonnegative) gives an integral optimal solution $y^*$, of value $\kappa$.

\emph{Reading off a set.} Let $A$ be the largest set of $\mathcal C_1$ with $y^*_{1,A}\ge1$, or $\varnothing$ if there is none, and $B$ likewise in $\mathcal C_2$. Each element $e$ is covered, so some $y^*_{k,C}\ge1$ with $e\in C$, and then
$e\in A\cup B$ because the sets of each chain are nested. So $E\setminus A\subseteq B$, and since ranks are nonnegative and monotone,
\[
\kappa=\sum_{k,C}r_k(C)\,y^*_{k,C}\ge r_1(A)+r_2(B)\ge r_1(A)+r_2(E\setminus A).
\]
With the easy inequality, this proves the theorem.
\end{proof}

\begin{remark}
The linear program has exponentially many constraints, so this proof is not an algorithm by itself. There are two ways out. Edmonds' combinatorial algorithm, with augmenting paths in an exchange graph, needs only a test of independence for each matroid, and
runs in polynomial time; the notebook of Tome~C1, Chapter~10, implements it. And the \emph{ellipsoid method} (Section~\ref{sec:lpspeed}) optimises over a polyhedron in polynomial time as soon as one can decide whether a point satisfies all the constraints, and if not exhibit one it violates;
for matroid polytopes this \emph{separation problem} is itself solvable in polynomial time. Gr\"otschel, Lov\'asz and Schrijver made this equivalence between separation and optimisation into a general theory in 1981.
\end{remark}

\section{How fast is the simplex method?}\label{sec:lpspeed}

Each pivot is cheap, but how many are needed? In practice, a small multiple of the number of constraints. In theory, the answer depends on the pivoting rule, and for every classical rule someone has built an example that forces exponentially many pivots. The first,
by Victor Klee and George Minty in 1972, is a squashed cube on which Dantzig's rule (enter the variable with the largest coefficient $\gamma_j$) visits all $2^n$ vertices:
\[
\begin{aligned}
\text{maximise }&\sum_{j=1}^n10^{n-j}x_j\\
\text{subject to }&\quad 2\sum_{j=1}^{i-1}10^{i-j}x_j+x_i\le100^{i-1}\ (1\le i\le n),\quad x\ge0.
\end{aligned}
\]
Figure~\ref{fig:kleeminty} shows the case $n=2$ (Exercise~\ref{ex:kleeminty}); the notebook counts $2^n-1$ pivots, and finds that Bland's rule, on the same cubes, needs fewer but still exponentially many, growing like the Fibonacci numbers (Avis and Chv\'atal showed in 1978 that Bland's rule can indeed be exponential).

\begin{figure}[ht]
\centering
\fig{ch07-kleeminty}
\caption{The Klee--Minty square: maximise $10x_1+x_2$ subject to $x_1\le1$, $20x_1+x_2\le100$, $x\ge0$ (vertical scale compressed). Dantzig's rule visits all four corners; entering $x_2$ first would reach the optimum $(0,100)$ in one pivot.}
\label{fig:kleeminty}
\end{figure}

Whether some pivoting rule is always polynomial is unknown. It would require, at least, that any two vertices of a polytope be joined by a path of polynomially many edges. In 1957 Warren Hirsch conjectured that $m-n$ edges always suffice for a polytope of dimension $n$
with $m$ facets; Francisco Santos disproved this in 2012 with a polytope of dimension $43$ and $86$ facets, but a polynomial bound remains open.

So the polynomial algorithms for linear programming came from elsewhere. In 1979 Leonid Khachiyan showed that the \emph{ellipsoid method}, invented for nonlinear optimisation by Shor, and by Yudin and Nemirovski, solves linear programs with rational data in time polynomial
in the number of bits of the input. The method is slow in practice, but its proof only needs a separation oracle, which is why it is a theoretical tool of great power, as in the remark of the previous section. In 1984 Narendra Karmarkar
found a polynomial \emph{interior point method}, walking through the inside of the polytope rather than along its edges, which competes with the simplex method in practice; these methods are the subject of Tome~C5. And in 2004 Daniel Spielman and Shang-Hua Teng
explained why the simplex method is fast in practice: if the data of any linear program are perturbed by a small random noise, the expected number of pivots of a version of the simplex method is polynomial (\emph{smoothed analysis}). Whether linear programs can be
solved in a number of arithmetic operations polynomial in $n$ and $m$ alone, independently of the size of the numbers, is still open; it is the ninth of the problems Stephen Smale proposed for the twenty-first century.

\begin{remark}
Linear programming had a forerunner. In 1939 the Soviet mathematician Leonid Kantorovich, consulted on how best to allocate work to the machines of a plywood trust, formulated such problems in his booklet \emph{Mathematical methods of organizing and planning production},
with a method of ``resolving multipliers'', the dual variables. It went unnoticed in the West until its translation in 1960. Kantorovich shared the 1975 Nobel prize in economics with Tjalling Koopmans for their contributions to the theory of the optimum allocation of resources.
\end{remark}

\begin{beauty}
A linear program comes with its shadow, the dual, and the best point of the one is a proof for the other: the last dictionary of the simplex method contains both the answer and the certificate. Integrality then reduces to determinants, and Edmonds' theorem
on two matroids comes out of the same mechanism after one uncrossing argument, with submodularity doing exactly the work that the triangle inequality does elsewhere. Behind it all is one geometric fact, that a point outside a convex set can be separated from it by a hyperplane.
\end{beauty}

\section{Exercises}

\begin{exercise}[$\star$]\label{ex:bakerydual}
Using the final dictionary of the bakery, find the best morning when there are $24+\delta$\,kg of flour, for which $\delta$ this stays valid, and the new profit. Compare with $y_1=\tfrac34$.
\end{exercise}

\begin{exercise}[$\star$]\label{ex:kleeminty}
Run the simplex method on the Klee--Minty square of Figure~\ref{fig:kleeminty} with Dantzig's rule, writing the three dictionaries, then with the opposite first choice.
\end{exercise}

\begin{exercise}[$\star$]\label{ex:degenerate}
Show that if no pivot is degenerate, the simplex method stops with any choice of entering and leaving variables.
\end{exercise}

\begin{exercise}[$\star$]\label{ex:extreme}
Let $P=\{x:Ax\le b\}$. Show that $v\in P$ is a vertex if and only if there are no distinct $u,w\in P$ with $v=\tfrac12(u+w)$.
\end{exercise}

\begin{exercise}[$\star\star$]\label{ex:bfsvertex}
Show that the vertices of $\{x\in\mathbb R^n:Ax\le b,\ x\ge0\}$ are exactly the restrictions to $x_1,\dots,x_n$ of the basic solutions of the feasible dictionaries.
\end{exercise}

\begin{exercise}[$\star$]\label{ex:dualrules}
Show that the dual of~\eqref{eq:dual}, written in standard form, is~\eqref{eq:primal}. Then show that the dual of ``maximise $c^{\mathsf T}x$ subject to $Ax\le b$'' with $x$ of any sign is ``minimise $b^{\mathsf T}y$ subject to $A^{\mathsf T}y=c$, $y\ge0$'', and that
an equation constraint in the primal gives a dual variable of any sign.
\end{exercise}

\begin{exercise}[$\star\star$]\label{ex:gamevalue}
Two players play a zero-sum game with payoff matrix $G\in\mathbb R^{m\times n}$: the row player picks a probability vector $p$ on the rows, the column player one, $q$, on the columns, and the row player wins $p^{\mathsf T}Gq$ on average. Prove von Neumann's minimax
theorem (1928): $\max_p\min_q p^{\mathsf T}Gq=\min_q\max_p p^{\mathsf T}Gq$.
\end{exercise}

\begin{exercise}[$\star$]\label{ex:tuinterval}
Show that a $0/1$ matrix in which the $1$s of each row are consecutive is TU.
\end{exercise}

\begin{exercise}[$\star$]\label{ex:oddcycle}
For the triangle, show that $x=(\tfrac12,\tfrac12,\tfrac12)$ is a vertex of $\{x\ge0:Mx\le\mathbf1\}$, where $M$ is the incidence matrix, and compute $\det M$. What does Edmonds' theorem on the matching polytope add to these constraints?
\end{exercise}

\begin{exercise}[$\star\star$]\label{ex:koniglp}
Deduce K\H onig's theorem from Theorem~\ref{thm:matroidintersection}, using the two partition matroids of a bipartite graph.
\end{exercise}

\section{Hints}

\begin{hint}{ex:bakerydual}
More flour means the slack $x_3$ can be $-\delta$. Put $x_3=-\delta$, $x_4=0$ in the final dictionary and check signs.
\end{hint}
\begin{hint}{ex:kleeminty}
Name the slacks $x_3$ and $x_4$. With Dantzig's rule, $x_1$ enters first because $10>1$.
\end{hint}
\begin{hint}{ex:degenerate}
Each pivot then increases $\zeta$ strictly, and $\zeta$ is determined by the basis.
\end{hint}
\begin{hint}{ex:extreme}
If the tight rows span, $a_i^{\mathsf T}u\le b_i$ and $a_i^{\mathsf T}w\le b_i$ average to an equality. Otherwise move along a direction orthogonal to the tight rows.
\end{hint}
\begin{hint}{ex:bfsvertex}
At a vertex, $n$ independent tight constraints are of the form ``a variable among $x_1,\dots,x_{n+m}$ is $0$''. Make these variables nonbasic.
\end{hint}
\begin{hint}{ex:dualrules}
Write equations as pairs of inequalities and free variables as differences, then dualise mechanically.
\end{hint}
\begin{hint}{ex:gamevalue}
For fixed $p$, the minimum over $q$ is attained at a column. Write the row player's problem as a linear program in $(p,v)$ and dualise.
\end{hint}
\begin{hint}{ex:tuinterval}
In a square submatrix, subtract each column from the next one, from the right.
\end{hint}
\begin{hint}{ex:oddcycle}
The three constraints $x_{ab}+x_{ac}\le1$ and so on are tight. For the last question, think of the odd set formed by the three vertices.
\end{hint}
\begin{hint}{ex:koniglp}
For the partition matroid of the left side, $r_1(A)$ is the number of left vertices met by $A$.
\end{hint}

\section{Solutions}

\begin{solution}{ex:bakerydual}
\idea The final dictionary describes the whole affine space $\mathcal S$; changing $b_1$ only shifts the slack $x_3$.
\writeup With $24+\delta$\,kg of flour, the slack becomes $x_3'=24+\delta-6x_1-4x_2=x_3+\delta$, where $x_3$ is the old expression. The final dictionary, written for $x_3=x_3'-\delta$, has the same coefficients and new constants:
with $x_3'=x_4=0$ it gives $x_1=3+\tfrac\delta4$, $x_2=\tfrac32-\tfrac\delta8$, $x_5=\tfrac52+\tfrac{3\delta}8$, $x_6=\tfrac12+\tfrac\delta8$, and $z=21+\tfrac34\delta-\tfrac34x_3'-\tfrac12x_4$. As long as these values are nonnegative, that is $-4\le\delta\le12$,
the dictionary is feasible and optimal, and the profit is $21+\tfrac34\delta$. For $\delta=1$: $(3.25,1.375)$ and $21.75$. Each extra kilogram is worth exactly $y_1=\tfrac34$ euro, within this range: the dual variable is the marginal value of the resource.
\end{solution}

\begin{solution}{ex:kleeminty}
\idea Follow the rule; the large coefficient of $x_1$ lures the method into the long way round.
\writeup Start with $x_3=1-x_1$, $x_4=100-20x_1-x_2$, $z=10x_1+x_2$. Dantzig's rule enters $x_1$; the ratios are $1$ and $5$, so $x_3$ leaves: $x_1=1-x_3$, $x_4=80+20x_3-x_2$, $z=10-10x_3+x_2$ (corner $(1,0)$). Now only $x_2$ can enter, $x_4$ leaves:
$x_2=80+20x_3-x_4$, $z=90+10x_3-x_4$ (corner $(1,80)$). Then $x_3$ enters, and only the row $x_1=1-x_3$ limits it: $x_3=1-x_1$, $x_2=100-20x_1-x_4$, $z=100-10x_1-x_4$, optimal at $(0,100)$ after three pivots. Entering $x_2$ first, only $x_4$
limits it, at $100$: $x_2=100-20x_1-x_4$, $z=100-10x_1-x_4$, optimal after one pivot.
\end{solution}

\begin{solution}{ex:degenerate}
\idea A strictly increasing quantity determined by the state cannot revisit a state.
\writeup If $\beta_\ell>0$, Proposition~\ref{prop:pivot} gives a strict increase of $\zeta$, since $\gamma_e>0$. If every pivot is nondegenerate, the values $\zeta$ of the successive dictionaries strictly increase, so no basis is visited twice
($\zeta$ is a function of the dictionary, hence of the basis by Lemma~\ref{lem:dictunique}). There are finitely many bases, so the method stops.
\end{solution}

\begin{solution}{ex:extreme}
\idea A vertex is pinned down by its tight constraints; a non-vertex can wiggle along a free direction.
\writeup If $v$ is a vertex and $v=\tfrac12(u+w)$ with $u,w\in P$, then for each tight $i$, $b_i=a_i^{\mathsf T}v=\tfrac12(a_i^{\mathsf T}u+a_i^{\mathsf T}w)$ with both terms $\le b_i$, so both are equal to $b_i$: $u$ and $w$ satisfy the tight equations, whose only solution is $v$,
so $u=w=v$. If $v$ is not a vertex, pick $d\ne0$ orthogonal to the tight rows; for small $\lambda>0$, $v\pm\lambda d$ satisfy the non-tight constraints (strict inequalities, by continuity) and the tight ones (with equality), and $v$ is their midpoint.
\end{solution}

\begin{solution}{ex:bfsvertex}
\idea The $n+m$ constraints of the standard form are exactly ``$x_k\ge0$'' for the $n+m$ variables, slacks included.
\writeup Identify $x\in\mathbb R^n$ with the point of $\mathcal S$ it determines; the constraint $x_j\ge0$ ($j\le n$) and the constraint $i$ ($x_{n+i}\ge0$) are tight when the corresponding variable is $0$, and each variable $x_k$ is an affine function
$\ell_k(x)$ of $x\in\mathbb R^n$. Given a feasible dictionary, the $n$ nonbasic variables vanish at its basic solution, and every $x_k$ is an affine function of $x_N$; in particular $x\mapsto x_N$ is an affine bijection of $\mathbb R^n$ onto $\mathbb R^N$, so the
linear parts of the $\ell_k$, $k\in N$, are independent, and the basic solution is a vertex. Conversely, at a vertex $v$, the tight constraints span, so we can choose $n$ of them with independent linear parts, $\ell_k(v)=0$ for $k\in N$, $|N|=n$. Then
$x\mapsto(\ell_k(x))_{k\in N}$ is an affine bijection, so every variable is an affine function of $x_N$: this is a dictionary with nonbasic set $N$, its basic solution is $v$ (the unique point with $x_N=0$), and it is feasible because $v\ge0$ in all
$n+m$ variables.
\end{solution}

\begin{solution}{ex:dualrules}
\idea Dualising is mechanical; only the bookkeeping of signs needs care.
\writeup The dual~\eqref{eq:dual} in standard form is: maximise $-b^{\mathsf T}y$ subject to $-A^{\mathsf T}y\le-c$, $y\ge0$. Its dual is: minimise $-c^{\mathsf T}x$ subject to $(-A^{\mathsf T})^{\mathsf T}x\ge-b$, $x\ge0$, that is, maximise $c^{\mathsf T}x$ subject to
$Ax\le b$, $x\ge0$. For a free $x=x^+-x^-$, the program ``maximise $c^{\mathsf T}x^+-c^{\mathsf T}x^-$ subject to $Ax^+-Ax^-\le b$, $x^\pm\ge0$'' has the dual ``minimise $b^{\mathsf T}y$ subject to $A^{\mathsf T}y\ge c$, $-A^{\mathsf T}y\ge-c$, $y\ge0$'', that is $A^{\mathsf T}y=c$.
Symmetrically, an equation $a^{\mathsf T}x=\beta$ is the pair $a^{\mathsf T}x\le\beta$, $-a^{\mathsf T}x\le-\beta$, with dual variables $y^+,y^-\ge0$ which only appear through $y^+-y^-$, a variable of any sign.
\end{solution}

\begin{solution}{ex:gamevalue}
\idea Each player's best guaranteed outcome is a linear program, and the two programs are dual.
\writeup For fixed $p$, $p^{\mathsf T}Gq$ is linear in $q$ on the simplex, so its minimum is attained at a vertex: $\min_qp^{\mathsf T}Gq=\min_j(G^{\mathsf T}p)_j$. So $\max_p\min_q$ is the value of the program: maximise $v$ subject to $v-(G^{\mathsf T}p)_j\le0$ for all $j$,
$\sum_ip_i=1$, $p\ge0$, $v$ free. By the rules of Exercise~\ref{ex:dualrules}, with dual variables $q_j\ge0$ for the first constraints and $w$ free for the equation, its dual is: minimise $w$ subject to $w-(Gq)_i\ge0$ for all $i$ (one constraint for each $p_i\ge0$),
$\sum_jq_j=1$ (for the free $v$), $q\ge0$; its value is $\min_q\max_i(Gq)_i=\min_q\max_pp^{\mathsf T}Gq$. Both programs are feasible and bounded (by the largest and smallest entries of $G$), so by Theorem~\ref{thm:fundamental} the first has an optimal solution, and by Theorem~\ref{thm:strongduality} the
values are equal.
\end{solution}

\begin{solution}{ex:tuinterval}
\idea Differences of consecutive columns turn intervals into a $+1$ and a $-1$: an incidence matrix of a digraph.
\writeup A square submatrix $Q$ of such a matrix keeps the property (the chosen columns, in order, still meet each row in a block). For $j=k,k-1,\dots,2$, replace column $j$ of $Q$ by column $j$ minus column $j-1$; the determinant does not change. A row whose $1$s
occupied columns $a$ to $b$ now has a $+1$ in column $a$, a $-1$ in column $b+1$ if $b<k$, and $0$ elsewhere. The transpose of the new matrix has at most a $+1$ and a $-1$ in each column, so it is TU by Theorem~\ref{thm:tucriterion} (all rows in one group),
and $\det Q\in\{0,\pm1\}$.
\end{solution}

\begin{solution}{ex:oddcycle}
\idea Odd cycles are exactly what breaks total unimodularity for matchings.
\writeup With vertices $a,b,c$ and edges $ab,ac,bc$, $M=\left(\begin{smallmatrix}1&1&0\\1&0&1\\0&1&1\end{smallmatrix}\right)$ and $\det M=-2$. At $x=(\tfrac12,\tfrac12,\tfrac12)$ the three vertex constraints are tight and $M$ is invertible, so $x$ is a vertex, of value $\tfrac32$,
while the largest matching has size $1$. Edmonds proved in 1965 that adding the constraints $x(E[U])\le\tfrac{|U|-1}2$ for all odd sets $U$ of vertices, where $E[U]$ is the set of edges inside $U$, gives exactly the convex hull of the matchings; here $U=\{a,b,c\}$ gives $x(E)\le1$, which cuts off $x$.
\end{solution}

\begin{solution}{ex:koniglp}
\idea The two rank functions count left and right endpoints, and the min--max formula becomes a vertex cover.
\writeup Let $G$ be bipartite with sides $L$ and $R$, and $E$ its edges. In $M_1$ a set of edges is independent if no two share a left endpoint, and $r_1(A)$ is the number of left vertices met by $A$; $M_2$ is the same on the right. Common independent sets
are matchings. For $A\subseteq E$, the left endpoints of $A$ and the right endpoints of $E\setminus A$ form a vertex cover of size at most $r_1(A)+r_2(E\setminus A)$. Conversely, for a vertex cover $K$, let $A$ be the set of edges with left endpoint in $K$; then every
edge of $E\setminus A$ has its right endpoint in $K$, so $r_1(A)+r_2(E\setminus A)\le|K\cap L|+|K\cap R|=|K|$. So the minimum in Theorem~\ref{thm:matroidintersection} is the smallest size of a vertex cover, and the maximum is the largest matching.
\end{solution}

\section*{Sources and further reading}

\begin{itemize}
\item \emph{The stories.} R.~W.~Cottle and M.~H.~Wright, \emph{Remembering George Dantzig}, SIAM News, April 2006, for the Neyman homework (``sometimes mistakenly considered an urban legend'') and the model of mid-1947; D.~J.~Albers and C.~Reid, \emph{An interview with George B.~Dantzig: the father of linear
programming}, College Mathematics Journal 17 (1986), 293--314, Dantzig's own account; S.~I.~Gass, \emph{George B.~Dantzig 1914--2005}, in \emph{Memorial Tributes} 12, National Academy of Engineering, 2008, for the simplex method at the end of the summer of 1947; G.~B.~Dantzig,
\emph{Linear programming}, Operations Research 50 (2002), 42--47, for the meeting with von Neumann in October 1947, as quoted by B.~von~Stengel, \emph{Zero-sum games and linear programming duality}, arXiv:2205.11196 (2023), which also gives the history of the equivalence of games and linear programs.
\item V.~Chv\'atal, \emph{Linear Programming}, W.~H.~Freeman, 1983, for dictionaries, the two-phase method and the proof of Bland's theorem given here; R.~G.~Bland, \emph{New finite pivoting rules for the simplex method}, Mathematics of Operations Research 2 (1977), 103--107; D.~Gale, H.~W.~Kuhn and A.~W.~Tucker,
\emph{Linear programming and the theory of games}, in \emph{Activity Analysis of Production and Allocation} (T.~C.~Koopmans, ed.), Wiley, 1951, 317--329; J.~Farkas, \emph{Theorie der einfachen Ungleichungen}, Journal f\"ur die reine und angewandte Mathematik 124 (1902), 1--27.
\item A.~J.~Hoffman and J.~B.~Kruskal, \emph{Integral boundary points of convex polyhedra}, and I.~Heller and C.~B.~Tompkins, \emph{An extension of a theorem of Dantzig's}, both in \emph{Linear Inequalities and Related Systems} (H.~W.~Kuhn and A.~W.~Tucker, eds.), Annals of Mathematics Studies 38, Princeton, 1956,
223--246 and 247--254; J.~Edmonds, \emph{Submodular functions, matroids, and certain polyhedra}, in \emph{Combinatorial Structures and their Applications} (Calgary, 1969), Gordon and Breach, 1970, 69--87; J.~Edmonds, \emph{Maximum matching and a polyhedron with $0,1$-vertices}, Journal of Research of the National Bureau of
Standards 69B (1965), 125--130. A.~Schrijver, \emph{Theory of Linear and Integer Programming}, Wiley, 1986, and \emph{Combinatorial Optimization}, Springer, 2003, Chapters~19 and~41, are the reference works.
\item D.~Avis and V.~Chv\'atal, \emph{Notes on Bland's pivoting rule}, Mathematical Programming Study 8 (1978), 24--34; V.~Klee and G.~J.~Minty, \emph{How good is the simplex algorithm?}, in \emph{Inequalities III} (O.~Shisha, ed.), Academic Press, 1972, 159--175; L.~G.~Khachiyan, \emph{A polynomial algorithm in linear programming}, Soviet Mathematics Doklady 20 (1979), 191--194; N.~Karmarkar, \emph{A new polynomial-time
algorithm for linear programming}, Combinatorica 4 (1984), 373--395; D.~A.~Spielman and S.-H.~Teng, \emph{Smoothed analysis of algorithms: why the simplex algorithm usually takes polynomial time}, Journal of the ACM 51 (2004), 385--463; F.~Santos, \emph{A counterexample to the Hirsch conjecture}, Annals of Mathematics 176 (2012),
383--412; S.~Smale, \emph{Mathematical problems for the next century}, Mathematical Intelligencer 20(2) (1998), 7--15; M.~Gr\"otschel, L.~Lov\'asz and A.~Schrijver, \emph{The ellipsoid method and its consequences in combinatorial optimization}, Combinatorica 1 (1981), 169--197.
\item L.~V.~Kantorovich, \emph{Mathematical methods of organizing and planning production} (1939), English translation in Management Science 6 (1960), 366--422.
\end{itemize}
